
\subsubsection{6.1 The Dual Finding: Existence Without Expected Mechanism}

The core finding is a dual result. The trace signal is real: per-sample last-fc Hessian trace achieves $AUROC\geq 0.85$ for minority membership detection in 4/5 seeds, with the signal being ERM-induced (epoch-0 control). The expected mechanism is absent: at the checkpoint of peak discriminability, minority training confidence is 0.97--1.0 --- indistinguishable from majority in terms of the confidence channel p\_i($1-p_{i})$.

This combination is more informative than either result alone. If only H-E3 had been tested, the confidence-differential mechanism would be the default attribution. If only H-M1 had been tested, there would be no evidence that the trace signal is discriminative at all. Together, they establish that the signal exists and identify the correct channel to investigate next.

\subsubsection{6.2 What Drives the Trace Signal? The Feature-Norm Channel}

For a linear cross-entropy head with weight W $\in \mathbb{R}^{C\times d}$ and penultimate feature x\_i $\in \mathbb{R}^d$:

$$\mathrm{Tr}(H_i^{fc}) \propto \|x_i\|^2 \cdot p_i(1 - p_i)$$

At t*, p\_i($1-p_{i})\rightarrow 0$ for all training samples (both minority and majority), as confirmed by H-M1. The confidence term is inactive. The only surviving candidate driver of the trace asymmetry is $\|x_{i}\|^2$ --- systematic differences in the L2 norm of the penultimate-layer feature representation between minority and majority samples.

LaBonte et al. [2024] provide a group-level prediction consistent with this interpretation: minority group covariance matrices have larger spectral norm than majority groups within the same class. Our finding extends this from the group-level covariance setting to the per-sample trace trajectory, and identifies $\|x_{i}\|^2$ as the specific channel that would need to show differential behavior.

\textbf{This interpretation is a candidate, not a verified claim.} Penultimate-layer norms $\|x_{i}\|^2$ have not been directly measured per sample, and it has not been confirmed that minority norms exceed majority norms. Verifying this --- by measuring feature norms at t* and correlating them with trace ranks --- is the primary direction for future mechanistic investigation.

\subsubsection{6.3 Relation to LaBonte and Muthukumar [2026] Phase I/II Theory}

LaBonte and Muthukumar [2026] prove for an XOR spurious feature model that SGD learns the spurious feature first (Phase I), then the core feature (Phase II). This predicts that during Phase I, minority samples remain near the decision boundary, keeping p\_minority(t)$\approx 0.5$ while p\_majority(t)>0.9.

The present results show that on ResNet-50 on Waterbirds, a transient version of this Phase I differential is present at epoch 1 but does not persist to t\textit{. ERM memorizes minority training samples --- driving their training confidence to $\geq 0.97$ --- before the trace ratio R(t) reaches its maximum. The mechanism operates on a different timescale than t}.

Two non-exclusive interpretations: (1) the Phase I window for ResNet-50 on Waterbirds closes quickly (by epoch 5 at the latest), so the confidence differential is genuinely transient; (2) ResNet-50 with ImageNet pretraining can memorize minority samples via the pretrained feature representation even without the spurious correlation, and this memorization outpaces the Phase I dynamics. Distinguishing these requires varying the spuriosity level and model capacity, which remains for future work. Additionally, the H-M1 test measured training-set confidence, where ERM optimization drives all samples toward high confidence; the mechanism may still operate on the held-out distribution, where minority samples encounter distribution shift.

\subsubsection{6.4 Limitations}

\textbf{Scope:} All experiments use ResNet-50 on Waterbirds v1.0 with 95\% spuriosity. Generalization to other architectures (ViT, ResNet-18), datasets (CelebA, MultiNLI), or spuriosity levels is not established.

\textbf{Mechanism unverified:} The feature-norm channel ($\|x_{i}\|^{2})$ is a candidate explanation, not a verified claim. The direct next step is to measure per-sample penultimate-layer norms at t* and test whether minority norms exceed majority norms at the same rank order as Hessian traces.

\textbf{No DFR application:} The trace signal has not been tested as a sample-weight or resampling guide for last-layer retraining (analogous to JTT upweighting). The existence result (AUROC>0.85) is necessary but not sufficient for practical WGA improvement --- the signal quality needs to translate through the DFR pipeline. Three downstream experiments (H-M2, H-M3, H-M4) were blocked when H-M1 failed its gate.

\textbf{Training-set confidence only:} The H-M1 mechanism test uses training-set confidence. The pre-registered gate was designed for training-set data, but future work should examine held-out confidence trajectories, where distribution shift may expose boundary-condition behavior that is masked by ERM memorization on the training set.

\textbf{Seed 1 non-monotone trajectory:} Seed 1 fails the H-E3 gate due to a non-monotone AUROC trajectory at t=$1\rightarrow 5$ (Spearman $\rho$=0.70, threshold 0.80). This may reflect stochastic optimization dynamics at a specific random initialization. Denser checkpointing (every epoch from t=0 to t=10) would clarify whether this is a genuine double-peak or a checkpoint-granularity artifact.

\textbf{Scalability:} Per-checkpoint trace computation takes approximately 8 minutes on an H100 for 4795 samples with K=50. Scaling to datasets with more than 100K samples or architectures with larger last-fc layers would require further engineering.

\subsubsection{6.5 Implications for Annotation-Free Minority Detection}

Per-sample last-fc Hessian trace is established as a candidate second-order annotation-free proxy for minority group membership --- the first such signal to achieve AUROC>0.85 on Waterbirds without group labels, with a verified epoch-0 control ruling out pretrained-artifact confounds.

Key practical questions remaining: (1) does the trace proxy translate to WGA improvement via DFR-style last-layer retraining; (2) does the signal generalize beyond Waterbirds; (3) what is the minimal computation needed (can K be reduced below 50, or can the epoch-0 control be replaced by a cheaper baseline check). If the feature-norm mechanism is confirmed, directly measuring $\|x_{i}\|^2$ as an annotation-free proxy might be computationally cheaper than the full Hutchinson computation.
